% Transcription — fonds Alexandre Grothendieck, shelfmark 87, batch 1 (pages 1-20).
% Established from the facsimile archives/batches/87/batch-01.pdf (a mirror of
% https://grothendieck.umontpellier.fr/87.pdf), per the skill
% transcribe-grothendieck. The folder is forty pages; this is the first of two
% batches.
%
% Sixteen of the twenty leaves carry his hand; fifteen are transcribed. Page 1
% is a tan wrapper sheet with nothing on it, and pages 16 and 20 are blank tan
% sheets of the same kind: all three are skipped, the gaps in the page numbers
% being the record. Page 15 is a tan cover carrying, in ink in the top right
% corner, the single word « Cubes », and in pencil, circled, « (13 p.) » — the
% latter is not obviously his and is not reproduced as his. The page takes no
% \page marker; « Cubes » becomes the heading of the second run and a \note
% says where it comes from. The inventory's title « Polyèdres convexes » is
% the archivists' (unbracketed); the first run's own heading, on page 2, is
% « Géométrie des convexes ». No pagination of his own on any leaf.
%
% Two runs. Pages 2-14 are a course-like development of the geometry of convex
% sets in a real affine space. Page 2 is a programme: convex parts, stability,
% barycentric description, the line case, the interior, and the radial
% homeomorphism between two compact convex bodies. Pages 3-5 define the
% facettes of a convex C as the classes of the relation R(x,y) « x = y or x
% and y interior to D_{x,y} n C », with E_F the affine span, the order F' < F
% iff F' c \bar F, finite sups (F_x = sup F_{x_i} for x a positive barycentre),
% extreme points, Krein-Milman in finite dimension, and finiteness of the
% facettes of a polyhedral part. Pages 5-6 take up polars (A° = {x' | <x,x'>
% >= -1 on A}), the bipolar theorem A°° = closed convex hull of A u {0}, and
% Hahn-Banach in finite dimension, proved on page 7 by reduction to the plane.
% Pages 7-10 prove that « secteurs polyédraux » (finite intersections of
% closed half-spaces = closed convex hulls of finitely many points and closed
% half-lines) are exchanged by polarity, via a lemma (page 9) that a finite
% intersection of half-spaces has finitely many extreme points. Pages 10-14:
% « Polarité et facettes » — a bijection F |-> F^0 between affine subspaces not
% through the origin, a theorem that polarity induces an order-reversing
% bijection between facettes of polar sectors, hence an anti-isomorphism
% Fac(C) = Fac(C') for polar polytopes; Corollary 3 that for a « polyèdre
% combinatoire strict » Phi, Phi° and the intervals Phi_{x,y} are again such;
% the dimension function (successor iff dimension goes up by one, Cor. 4-6);
% and Cor. 7, the « diamond » property: an interval of length two has exactly
% two middle elements, because a one-dimensional polytope is a segment. The
% notion « polyèdre combinatoire strict » is used on pages 12 and 14 without a
% definition on these leaves.
%
% Pages 17-19 (under « Cubes »): an intrinsic characterisation of an n-cube by
% its set Phi of affine subspaces (the affine spans of its faces), through
% axioms a)-e) — Phi contains the empty set and E, every element of Phi** is
% an intersection of maximal ones; each maximal F has a unique « opposite » F'
% with Inf(F,F') = empty, an involution sigma without fixed point on the set B
% of maximal elements; Phi \ {empty} -> Gamma_part(B/I) bijective; a margin
% says a)-c) mean Phi is a « cube combinatoire ». Then the reconstruction: the
% directions of the hyperplanes split the translation space T as a product of
% lines T_i, E as a product of T_i-torsors E_i, each carrying a two-point B_i,
% with T_i = Ker(k^{B_i} -> k), E_i = eps^{-1}(1), T = Ker(k^B -tr-> k^I). It
% breaks off at the foot of page 19 (« Si J c I, on considère »), and
% presumably continues in batch 2.
%
% Readings to check first: the abbreviation read as Fac(C) throughout (« F..(C) »
% on the leaves); the superscript on pages 10-11 read as ^0 (for the affine
% correspondence), distinct from the polar °; the superscript read as ** on
% Phi** (page 17); the letter read as \mathbb{H} on page 18; the word opening
% « Dém. » on page 7 (twice, \ill); the cancelled block on page 8 (a) => b),
% given as one \struck{\ill{}}); page 13 « dim C = 0 ou -1 ».
%
% Pass: Opus 5.5 (claude-opus-5-5), 2026-09-24 — first pass, unchecked against the pages by a human.

\documentclass[11pt,a4paper]{article}
% Le préambule commun à toutes les transcriptions du fonds Grothendieck.
%
% Il tient en une idée : **l'incertitude se compose, elle ne se résout pas.**
% Une transcription qui lisse un mot douteux a détruit la seule chose pour
% laquelle on la faisait — un lecteur doit pouvoir savoir, à chaque endroit,
% ce qui était sur le feuillet et ce qui a été ajouté. Les six macros qui
% suivent sont donc l'appareil critique, et non de la décoration.
%
% Les mêmes commandes sont reconnues par `scripts/render.mjs`, qui produit la
% vue de lecture du site. Une macro ajoutée ici doit l'être aussi là-bas,
% faute de quoi le rendu échouera bruyamment — ce qui est voulu : un rendu
% silencieusement faux est pire qu'un rendu absent.

\ProvidesPackage{grothendieck}[2026/08/08 Transcriptions du fonds Grothendieck]

\usepackage{amsmath,amssymb,amsthm}
\usepackage{mathrsfs}
\usepackage{stmaryrd}
% Les diagrammes commutatifs sont partout dans ce fonds : c'est la langue
% même de la géométrie algébrique de Grothendieck, et une transcription qui
% les rendrait en prose ne transcrirait rien.
\usepackage{tikz-cd}
% Français par défaut — c'est la langue des feuillets — mais l'anglais est
% chargé aussi : la traduction et le résumé sont des documents anglais, et la
% typographie française (espaces avant les deux-points) y serait une faute.
% Ces documents basculent par \selectlanguage{english} en tête de corps.
\usepackage[english,french]{babel}
\usepackage{fontspec}
\usepackage{geometry}
\geometry{a4paper,margin=25mm,marginparwidth=28mm}
\usepackage{marginnote}
\usepackage{xcolor}
% ulem, not soul. `soul`'s \st and \ul are built on a character-by-character
% scan that XeLaTeX's font machinery breaks: the document fails with an
% undefined control sequence rather than degrading. ulem does the same two jobs
% and is Unicode-engine safe. `normalem` stops it hijacking \emph, which the
% transcriptions use for Grothendieck's own emphasis.
\usepackage[normalem]{ulem}
\usepackage{hyperref}
\hypersetup{colorlinks=true,linkcolor=black,urlcolor=black,pdfborder={0 0 0}}

% --- Métadonnées du lot -----------------------------------------------------
% Reprises telles quelles par le script de rendu, qui les affiche en tête de
% la vue de lecture. Elles fixent ce que le fichier prétend couvrir : sans
% elles, une transcription flotte sans point d'attache dans un fonds de seize
% mille feuillets.

\newcommand{\folder}[1]{\gdef\grothFolder{#1}}
\newcommand{\batch}[1]{\gdef\grothBatch{#1}}
\newcommand{\leaves}[2]{\gdef\grothFirst{#1}\gdef\grothLast{#2}}
\newcommand{\foldertitle}[1]{\gdef\grothFoldertitle{#1}}
\gdef\grothFolder{?}\gdef\grothBatch{?}\gdef\grothFirst{?}\gdef\grothLast{?}\gdef\grothFoldertitle{}

% --- Le résumé --------------------------------------------------------------
%
% Il ouvre la lecture modernisée, sous le titre « Résumé », et il est composé
% en petit. Ce n'est pas un ornement : c'est le seul passage du document qui
% ne soit pas des mathématiques, et un lecteur doit pouvoir le reconnaître
% avant de l'avoir lu — puis le sauter s'il connaît déjà le sujet, ou s'y
% arrêter s'il ne le connaît pas. Le filet qui le suit marque la reprise du
% corps.

\newenvironment{resume}
  {\small\setlength{\parskip}{0.45em}}
  {\par\medskip\hrule\medskip}

% --- Mots-clés ---------------------------------------------------------------
%
% En anglais, en clôture du résumé de la lecture modernisée : le vocabulaire
% moderne sous lequel chercher ce que les feuillets construisent. Le manifeste
% du site (`npm run manifest`) les extrait pour étiqueter la cote entière —
% cette ligne est la source unique des tags, il n'y a pas de fichier à côté.
\newcommand{\keywords}[1]{\par\smallskip\noindent
  {\footnotesize\textsc{Keywords} --- \textit{#1}}\par}

% --- L'appareil critique ----------------------------------------------------

% Le numéro de feuillet, en marge. C'est l'ancre qui permet de régler un
% désaccord sur une formule en regardant *un* feuillet plutôt qu'une chemise
% de six cents. Le site s'en sert aussi pour tourner les pages du fac-similé
% quand on fait défiler la transcription.
\newcommand{\leaf}[1]{%
  \marginnote{\footnotesize\color{gray!70}\textsf{#1}}%
  \label{leaf:#1}\ignorespaces}

% Illisible : on ne devine pas. Un mot inventé qui se lit comme les autres est
% le pire résultat possible.
\newcommand{\ill}{\textcolor{red!60!black}{[\,\ldots\,]}}

% Lecture douteuse : le mot est donné, et signalé comme incertain.
\newcommand{\uncertain}[1]{\uline{#1}}

% Ajout de l'éditeur — une lettre manquante, un mot sous-entendu.
\newcommand{\add}[1]{\textcolor{blue!50!black}{[#1]}}

% Biffé par Grothendieck lui-même. Ce qu'il a rayé fait partie du document :
% une rature dit souvent où la pensée a changé de direction.
\newcommand{\struck}[1]{\sout{#1}}

% Note du transcripteur, sur l'état matériel du feuillet.
\newcommand{\note}[1]{\par\smallskip{\small\color{gray!60!black}\textit{[#1]}}\par\smallskip}

% Note marginale de Grothendieck, distincte du corps du texte.
\newcommand{\marginal}[1]{\par\smallskip\noindent\hspace{1em}%
  {\small\color{gray!60!black}#1}\par\smallskip}

% --- Titre ------------------------------------------------------------------

\AtBeginDocument{%
  \begin{center}
    {\large\bfseries Fonds Alexandre Grothendieck --- cote n\textsuperscript{o}~\grothFolder}\\[2pt]
    {\itshape\grothFoldertitle}\\[4pt]
    {\small feuillets \grothFirst--\grothLast\ (lot \grothBatch)}\\[2pt]
    {\footnotesize\color{gray!60!black}Transcription établie d'après le fac-similé numérisé
     par l'Université de Montpellier.\\
     Les lectures incertaines sont soulignées, les illisibles notées [\,\ldots\,],
     les ajouts entre crochets.}
  \end{center}
  \vspace{1em}}

\endinput


\folder{87}
\batch{1}
\pages{1}{20}
\dating{s.d. — le groupe « Géométrie et topologie combinatoire » (69 à 102) est daté 1976-[vers 1986]}
\watermark{Édition de démonstration}
\foldertitle{Polyèdres convexes : notes manuscrites (s.d.)}

\begin{document}
\section{Géométrie des convexes}

\page{2}

\note{le titre est de sa main, souligné, en tête de la page 2. Ses rubriques
ouvertes par un astérisque sont rendues ici en liste sur cette page, puis en
sous-sections à partir de la page 3}

\begin{itemize}
  \item Partie convexe d'un espace affine réel : définition, stabilité par
    intersection, partie convexe \emph{engendrée} --- description en termes de
    combinaisons barycentriques, en termes de simplexes, comparaison avec
    espace affine \emph{engendré}. Enveloppe convexe d'\uncertain{une}
    \struck{\uncertain{réun.}} famille d'ens. convexes. Stabilité par images
    inv. et directes par appl. affines, par sommes.
  \item Cas de \struck{$\mathbf{R}$} droite : parties convexes sont les
    \emph{intervalles}. \struck{Une} partie \struck{convexe} $C$ d'un
    \struck{espace} espace affine réel $E$ est convexe ssi son intersection
    avec toute droite est un intervalle de ladite.

    Dimension d'un convexe.
  \item Intérieur d'un convexe (engendrant $E$ affinement) \uncertain{i.e.} de
    dim $E$
\end{itemize}

\begin{itemize}
      \item[a)] $\mathring{C} \neq \emptyset$
      \item[b)] Si $x \in \mathring{C}$, alors
        \[
          \mathring{C} = \bigcap_{\substack{D \text{ droite de } E \\ \text{passant par } x}} \operatorname{int}_D (D \cap C)
        \]
        \[
          \dot{C} = \bigcap_D \operatorname{bnd}_D (D \cap C)
        \]
        \[
          \overline{C} = \bigcap_D \overline{D \cap C}
        \]
        \note{à droite de ces formules, un croquis : une région $C$ hachurée
        par deux familles de droites parallèles. Le signe devant
        $\operatorname{bnd}_D$ est bien une intersection sur la page, là où
        l'on attendrait une réunion}
      \item[c)] Si $C$, $C'$ convexes compacts engendrant $E$, $x \in
        \mathring{C}$, $x' \in \mathring{C}'$, $\exists$ homéom. unique
        $h : C \simeq C'$ avec $h(x) = x'$, $h$ étant \struck{linéaire}
        affine sur les demi-droites issues de $x$ et les transformant en
        demi-droites \uncertain{parallèles}.
\end{itemize}

\page{3}

\subsection{Facettes d'un convexe}

$R(x,y)$ : ($x = y$ ou $x$ et $y$ sont intérieurs à $D_{x,y} \cap C$) est une
relation d'équivalence. Une partie $F$ de $C$ est dite une \emph{facette} si
$\forall x \in F$, on a $F = \{ y \in C \mid R(x,y) \}$
--- \struck{\uncertain{on exclut}} \uncertain{donc} les $F \neq \emptyset$
\struck{\uncertain{facettes}} --- $F$ est une classe suivant $R$. Notations
$E_C$, $E_F$\note{sous-espaces affines engendrés ; l'indice du premier $E$
est un $C$ appuyé. L'ensemble des facettes est noté par une abréviation lue
$\mathrm{Fac}(C)$ dans tout le lot}.

\emph{Point extrémal} $x$ : \struck{facette réduite à} $\{x\}$ est facette
$\Leftrightarrow$ $\forall D$ droite passant par $x$, $x$ n'est pas intérieur
à $D \cap C$ dans $D$.

(1) Supposons $C$ de dim finie $d$. $\exists !$ facette de dim $d$, c'est
$\operatorname{int}_{E_C}(C)$.

(2) Soit $F \in \mathrm{Fac}(C)$. \struck{Alors} $E_F \cap C$ est une réunion
de facettes de $C$, et \struck{\ill{}} les facettes de $C$ contenues dans
$E_F \cap C$ sont les facettes de $E_F \cap C$. Si $E$ de dim finie, on a
$E_F \cap C = \overline{F} \cap C$ ($= \overline{F}$ si $C$ fermé), et
$F = \operatorname{int}_{E_F}(E_F \cap C)$, \struck{Pour deux}

(3) Pour deux facettes $F$, $F'$, conditions équivalentes :
\[
  \begin{array}{ll}
    E_{F'} \subset E_F & \\
    E_{F'} \cap C \subset E_F \cap C & (\text{ou } E_{F'} \cap C \subset E_F) \\
    F' \subset E_F \cap C & (\text{ou } F' \subset E_F) \\
    F' \subset \overline{F} \cap C & (\text{ou } F' \subset \overline{F}) \\
    \overline{F'} \subset \overline{F} &
  \end{array}
\]
\note{une accolade réunit les deux dernières lignes, avec « si $E$ de dim
finie »}

C'est là une \emph{relation d'ordre} dans l'ens. $\mathrm{Fac}(C)$ des
facettes de $C$. Si $x, y \in C$, \struck{\ill{}}
\[
  F_x \prec F_y \Leftrightarrow
  \bigl[ x = y \text{ ou } x \neq y \text{ et }
  y \in \operatorname{Int}_{D_{xy}} (D_{xy} \cap C) \bigr]
\]
\note{un mot est noirci devant $F_x$}

\page{4}

(4) Soient $(x_i)_{i \in I}$ une famille finie de pts de $C$,
$(\lambda_i)_{i \in I}$ avec $\lambda_i > 0$, $\sum \lambda_i = 1$, donc
$x = \sum \lambda_i x_i \in C$. Ceci dit, on a
\[
  F_x = \sup F_{x_i}.
\]
(Par suite, dans $\mathrm{Fac}(C)$ les sup finis existent --- donc les sup
quelconques \uncertain{aussi} \uncertain{pourvu} que $\dim C < +\infty$).

(5) Soit $C' \subset C$ une partie de $C$ réunion de facettes, et soit
$\Phi \subset \mathrm{Fac}(C)$ défini par
\[
  \Phi = \{ F \in \mathrm{Fac}(C) \mid F \subset C' \}
\]
Pour que $C'$ soit convexe, il faut et suffit que $\Phi$ soit stable par sup
finis ; pour que $C'$ soit \struck{convexe et} fermé \uncertain{relat.} dans
$C$ (car dim $E < \infty$) il faut et suffit que $\Phi$ \struck{soit stable
par sup finis et} qu'avec une facette il contienne celles qui sont
\uncertain{majorées}. \struck{(5)} (Dim finie), (6) Pour que $C'$ soit de la
forme $\overline{F} \cap C$ ($F \in \mathrm{Fac}(C)$) il faut et suffit que
$\Phi$ soit fermé, convexe et \struck{soit} \struck{\ill{}}
\uncertain{réunion} de facettes. L'ens. des facettes \uncertain{relat.}
\uncertain{fermées} de $C$ \struck{\ill{}} est stable par intersections
quelconques (donc il y a des inf quelconques dans l'ens. des facettes de $C$).

\page{5}

(7) Supposons $C$ \uncertain{de} dim finie compact. Alors $C$ est l'enveloppe
convexe de l'ens. de ses pts extrémaux.

(8) Supposons $C$ partie polyédrale dans $E$ de dim finie (i.e. $C$ réunion
finie de simplexes $\Leftrightarrow$ réunion finie d'enveloppes convexes de
parties finies de $E$). Alors l'ens. des facettes de $C$ est fini.

\emph{Cor.} Polyèdres convexes : parties convexes compactes telles que
$\mathrm{Fac}(C)$ (ou $\mathrm{Fac}_0(C)$, ens. des pts extrémaux) soit fini.
\struck{Polaires}

\subsection{Polaires}

$E$, $E'$ espaces vectoriels duals l'un de l'autre. Pour partie $A$ de $E$,
on pose
\[
  A^\circ = \{ x' \in E' \mid \langle x, x' \rangle \geq -1 \ \forall x \in A \}
  = \bigcap_{x \in A} H(x)
  \qquad (\text{où } H(x) = \{ x' \in E' \mid \langle x, x' \rangle \geq -1 \})
\]
De même on définit $A'^\circ \subset E$ pour $A' \subset E'$.

\emph{Th. 1} Pour $A \subset E$, … $A^{\circ\circ}$ = env. conv. fermée
de $A \cup \{0\}$.

\emph{Corollaire} $A \mapsto A^\circ$ et $A' \mapsto A'^\circ$ sont des appl.
bijectives inv. l'une de l'autre entre l'ens. des parties convexes fermées de
$E$ contenant $0$, et l'ens. analogue

\page{6}

pour $E'$. \struck{En} De plus, pour de telles parties
\[
  A_1 \subset A_2 \Longleftrightarrow A_2^\circ \subset A_1^\circ
\]
(i.e. la bijection précédente renverse l'ordre)

\emph{Cor}
\[
  \overline{\mathrm{Env}}(A_1, A_2) = A_1^\circ \cap A_2^\circ
\]
\[
  (A_1 \cap A_2)^\circ = \overline{\mathrm{Env}}(A_1, A_2)
\]
\note{ainsi sur la page : aucun signe de polaire sur le membre de gauche de
la première formule ni sur les arguments du second membre de la seconde}

\emph{Cor} \struck{$A$} $0 \in \operatorname{int}(A) \Longleftrightarrow
A^\circ$ compact

$A$ \uncertain{cpct} $\Longleftrightarrow 0 \in$ \struck{\ill{}}
$\operatorname{int}(A^\circ)$

Le théorème est conséquence \struck{du théo.} \ill{} est équivalent :)
\note{la ligne est marquée à gauche d'un grand crochet}

\emph{Théorème 1'} Soit $C$ une partie de $E$. Pour que $C$ soit convexe et
fermé, il faut et suffit que $C$ soit intersection de demi-espaces fermés.

Résulte du

\emph{Théorème 1''} (Hahn-Banach) Soit $U$ une partie \emph{convexe ouverte}
de l'espace vectoriel réel $E$ de dim finie, $F \subsetneq E$ un sous-espace
vectoriel de $E$ tel que $F \cap U = \emptyset$, alors $\exists$ hyperplan $H$
de $E$ contenant $F$, tel que $H \cap U = \emptyset$.

\page{7}

\struck{Résultats} \emph{Dém.} \ill{} : remplaçant $U$ par
$\bigcup_{\lambda > 0} \lambda U$, OPS $U$ un \emph{cône} (i.e. stable par
multiplication par $\lambda > 0$) Il suffit de prouver que si $F$ n'est pas
un hyperplan, $\exists F' \supset F$, $F$ de codim 1 dans $F'$ i.e.
$\dim F' = \dim F + 1$. OPS $\dim E = \dim F + 2$. \ill{} : passant au
quotient par $F$, on se ramène au cas $\dim E = 2$, $F = 0$. OPS
$U \neq \emptyset$.\note{le même mot, non lu, ouvre les deux phrases}

On montre alors que ou bien $U$ est un demi-espace ouvert, ou bien $\exists$
base $e_1, e_2$ telle que
\[
  U = \{ \lambda_1 e_1 + \lambda_2 e_2 \mid \lambda_1, \lambda_2 > 0 \}
\]
\note{dans la marge gauche, deux axes perpendiculaires issus d'un point}

Dans les deux cas on gagne …

\subsection{Secteurs polyédraux}

\emph{Prop.} Soit $C$ partie de $E$ (espace affine de dim finie) Conditions
équivalentes :
\begin{itemize}
  \item[a)] $C$ intersection finie de demi-espaces fermés
  \item[b)] $C$ enveloppe convexe fermée d'un ens. fini de pts et de
    demi-droites fermées (qu'on peut supposer \struck{\uncertain{partant}}
    \uncertain{issues} \uncertain{d'une origine} un pt $x \in C$ donné à
    l'avance …)
  \item[b')] (si compact) $C$ est un polyèdre convexe
\end{itemize}

\page{8}

\emph{Dém.} Si $C = \emptyset$ les deux conditions sont satisfaites. Donc OPS
$C \neq \emptyset$. \struck{Soit $x \in C$, on le prend comme origine} OPS
$C$ convexe fermé, et \struck{OPS $C = C^{\circ\circ}$, \ill{} $C' =
C^\circ$ \ill{}} \uncertain{$C$ engendre $E$} (car on voit que les deux
conditions sur $C$ \uncertain{ne changent pas} si on remplace $E$ par $E_C$).
Soit $x \in \mathring{C}$, on le prend comme origine, on aura
\[
  C = C'^\circ, \quad C' = C^\circ.
\]
D'où $C'$ est compact.

\struck{\ill{}}\note{suit un long passage, encadré et barré de traits
obliques, qui commence par « a) $\Rightarrow$ b) » et où se lisent
« convexe fermée d'un ens. fini de pts et de demi-droites fermées », « b)
$\Rightarrow$ a) », « enveloppe convexe », « d'un ens. fini de pts »,
« intersection d'un ens. fini de demi-espaces fermés » et « on gagne par
polarité » ; il est remplacé par ce qui suit et par la page 9}

b) $\Rightarrow$ a) Si b) est vrai pour $C$, alors par polarité $C'$ est
int. finie de $\frac{1}{2}$-espaces fermés, donc (lemme ci-dessous) étant
compact, c'est un polyèdre convexe, i.e. env. conv. \struck{finie de} d'un
ens. fini de $E'$, d'où a) par polarité.

\page{9}

a) $\Rightarrow$ b) On sait (par polarité) que $C'$ est \struck{\ill{}} env.
conv. fermée d'un ens. fini de pts, donc d'après ce qui précède\struck{,
\ill{}} il est intersection finie de demi-espaces fermés\struck{\ill{}}
donc par polarité on trouve b) pour $E$, sous forme précisée.
\note{dans la marge gauche, un croquis : une bande entre deux droites
parallèles, un point marqué sur l'une, et une ligne brisée en zigzag entre
les deux}

\emph{Lemme} Soit $C$ \uncertain{compact} intersection finie de
demi-espaces fermés. Alors \struck{$C$ est un polyèdre convexe}
l'ens. des pts extrémaux de $C$ est fini (donc $C$ est un polyèdre convexe
s'il est compact). \struck{En d'autres termes, $C$ \ill{} est \ill{} fini de
pts extrémaux.}

Récurrence sur $\dim E$ (trivial si $\dim E = 0$) puis sur le
\struck{\ill{}} nb de demi-espaces \uncertain{qui entrent}. Trivial si
$\nu = 0, 1$ \struck{(\ill{} $n = 1$) Supposons}\note{le $\nu$ renvoie à une
ligne barrée, « sur $\nu = n + \dim E$ »} On est ramené au cas d'un convexe
$C_1$ n'ayant qu'un nb fini de pts extrémaux, et où on a
\[
  C = C_1 \cap H
  \qquad (H \text{ demi-espace fermé, limité par l'hyperplan } H_0)
\]
Soit $x \in C$ \struck{un pt extrémal}. Si $x \notin H_0$, alors $x$ est pt
extrémal de $C$ ssi il est pt extrémal de $C_1$. Si $x \in H_0$,
\struck{alors} $x$ est pt extrémal de $C$, ssi il est pt extrémal de
\struck{$C_1$ \ill{}} $H_0 \cap C$. On gagne.

\page{10}

\emph{Définition} \emph{Secteur polyédral} : partie satisfaisant aux
conditions équivalentes précédentes.

\emph{Prop.} $E$, $E'$ vectoriels de dim finie en dualité. Alors par
$A \mapsto A^\circ$, $A' \mapsto A'^\circ$ on trouve des correspondances 1-1
inverses l'une de l'autre entre secteurs polyédraux de $E$
\uncertain{contenant} l'origine, et sect. pol. de $E'$ contenant l'origine.

\subsection{Polarité et facettes}

$E$, $E'$ vectoriels de dim finie en dualité. Correspondance biunivoque
entre sous-espaces affines $F$ de $E$ ne contenant pas l'origine, et
sous-espaces affines $F'$ de $E'$ ne contenant pas l'origine, avec
\[
  \dim F + \dim F^{0} = n - 1 \qquad (n = \dim E = \dim E')
\]
\marginal{Notations $F^{0}$,}\note{l'exposant de $F^{0}$ est un petit rond
ouvert, tracé autrement que le $^\circ$ des polaires ; il est rendu $^{0}$
ici et page 11}

Correspondance caractérisée par les propriétés :
\begin{itemize}
  \item[a)] Renverse la relation d'incidence
  \item[b)] À $x \in E - \{0\}$ correspond
    $\{x\}^{0} = H(x)^{\bullet} = (\{x\}^\circ)^{\bullet}$
    \note{le point en exposant désigne sans doute le bord, comme
    $\dot{C}$ page 2}
\end{itemize}

\page{11}

\emph{Terminologie} Variété d'appui d'un convexe $C$ = variété affine de la
forme $E_F$. Les variétés d'appui sont en corr. 1-1 avec les facettes,
relation d'ordre \uncertain{et} relation d'incidence

\emph{Th} Soit $C$, $C'$ deux secteurs polyédraux polaires l'un de l'autre de
$E$, $E'$ resp.t Alors il existe une unique \struck{\ill{}} application
$\varphi$ de l'ens. des facettes $F$ de $C$ telles que $0 \notin
\overline{F}$ (i.e. $0 \in E_F$) dans l'ens. des facettes $F'$ de $C'$ telles
que $0 \notin \overline{F'}$ (i.e. $0' \in E_{F'}$), caractérisée par la
propriété que \struck{si $F' = \varphi(F)$}
\[
  E_{\varphi(F) = F'} = (E_F)^{0}
\]
\note{les deux « i.e. » portent $\in$ et non $\notin$ sur la page. Devant
la formule, un mot et un $E$ sont barrés}

Cette application est bijective, et renverse la relation d'ordre entre
facettes

\emph{Cor. 1} L'ens. des facettes d'un secteur polyédral est fini (prendre
l'origine \emph{intérieure} au secteur polyédral, et appliquer dualité)

\emph{Cor. 2} Supposons $C$, $C'$ cpts \struck{\ill{}} i.e. \struck{$E$}
$0_E$, $0_{E'}$ intérieurs à $C$, $C'$ resp.t, i.e. $C$, $C'$ des polyèdres.
Alors $\varphi : \mathrm{Fac}(C) \simeq \mathrm{Fac}(C')$

\page{12}

(antiisomorphisme d'ens. ordonnés).

\emph{Corollaire 3} Soit $\Phi$ un polyèdre combinatoire strict. Alors
\begin{itemize}
  \item[a)] $\Phi^\circ$ est un pol. combinatoire strict
  \item[b)] Si $x, y \in \Phi$, avec $x \leq y$, alors l'ens. ordonné
    $\Phi_{x,y} = \{ z \in \Phi \mid x \leq z \leq y \}$ est un polyèdre
    combinatoire strict.
\end{itemize}
\note{la notion de polyèdre combinatoire strict n'est pas définie sur ces
feuillets}

\emph{Dém.} a) Résulte de Cor. 2. On sait que $\Phi_{0,x}$ est un polyèdre
combinatoire strict, ($0$ plus petit élt) donc par dualité aussi $\Phi_{x,1}$
($1$ plus grand él.) $= (\Phi^\circ_{0,x})^\circ$. Appliquant ce dernier
résultat à $x \in \Phi_{0,y}$, on trouve que $\Phi_{xy}$ est un polyèdre
combinatoire strict.

\struck{Cor. 4 Soient $x, y$ deux facettes}

\struck{Cor. 4 Soit $F$ une facette d'un p. $C$ un polyèdre convexe, $F$ une
facette dudit, $i \in \mathbf{Z}$ tel que $-1 \leq i \leq d = \dim F$. Alors
$\exists$ facette de $F'$ incidente à $F$ telle que $\dim F' =$}

\page{13}

\struck{Si $d = -1$ ou $d = 0$, c'est évident. Supposons $d \geq 1$, i.e.
$F \neq \emptyset$, on sait que $\overline{F}$ est l'enveloppe convexe de
l'ens. de ses pts extrémaux, donc $\exists F' \in \mathrm{Fac}(C)$ avec
$\dim F' = 0$, et $F' \leq F$, \ill{} $F' < F$.}\note{ce passage et le
corollaire barré de la page 12 sont encadrés et barrés de traits obliques ;
l'énoncé revient page 14 comme Cor. 6}

\subsection{Application à la fonction dimension}

\emph{Cor. 4} Soient $C$ un polyèdre convexe, $F_1, F_2 \in \mathrm{Fac}(C)$
avec $F_1 < F_2$. Pour que $F_2$ soit un \emph{successeur} de $F_1$ (i.e.
qu'il n'existe pas $F$ tel que $F_1 < F < F_2$) il faut et suffit qu'on ait
\[
  \dim F_2 = \dim F_1 + 1
\]

Comme $\dim F$ est fonction \emph{strictement} croissante sur
$\mathrm{Fac}(C)$, la suffisance est claire. Pour la nécessité,
\struck{\ill{} dans \ill{} $\dim F$ \ill{}} OPS (\uncertain{dualité})
$F_2 =$ \struck{$\mathring{C}$} $1$ (plus grande facette), \struck{et $F_2$}
puis en passant aux polaires (échangeant les rôles de $F_1$ et $F_2$) que
$F_1 =$ \struck{$\emptyset$} $0$ (plus petite facette), et :
\uncertain{à nouveau} $F_2 = 1$. Donc l'assertion devient maintenant : Pour
qu'il n'existe aucune facette distincte de $\mathring{C}$ et de la facette
vide, il faut et suffit que $\dim C \uncertain{=} 0$ ou $-1$ i.e. que $C$
\uncertain{vide} \struck{\uncertain{réduit}} à un pt. \struck{\ill{}}
(\uncertain{trivial}, en effet, on prend pts extrémaux !)

\emph{Cor. 5} La fonction $d : F \mapsto \dim F$ sur $\mathrm{Fac}(C)$ (du
cor. 4) est caractérisée par les conditions :
\begin{itemize}
  \item[a)] $d(F_2) = d(F_1) + 1$ si $F_2$ est un successeur de $F_1$
  \item[b)] $d(\emptyset) = -1$ ($\emptyset$ plus petit él.). De plus $d$ est
    str. croissant
\end{itemize}

\page{14}

\emph{Cor. 6} Soit $F \in \mathrm{Fac}(C)$ ($C$ polyèdre convexe) Alors pour
tout $-1 \leq i \leq d = \dim F$, existe $F' \leq F$ telle que
$\dim F' = i$.

\struck{Clair si $d = -1$ ou $0$, sinon \ill{}} \struck{Récurrence sur $d$,
clair si $d = -1, 0$. Si $d > 0$, $\exists F'$ tel que $F' \neq \emptyset,
F$, $\emptyset < F' < F$ (cor. 4) Donc $\dim F' = d' < d$. Si $i \leq d'$
on gagne par hyp. de récurrence appliquée à $F'$ \ill{} sinon on gagne par
l'hyp. de récurrence}\note{bloc encadré et barré de traits obliques}

\uncertain{Formel} à partir du cor. 5, qui implique que si $x \leq y$,
$x, y \in \mathrm{Fac}(C)$, alors $\exists$ chaîne
$x_0 = x < x_1 < \cdots < x_n = y$, avec $x_{i+1}$ succ. de $x_i$
($0 \leq i \leq n-1$), et qu'alors \struck{$d(y) - d(x) = n$}
\struck{\ill{}} $d(x_i) - d(x_0) = i$ i.e. $d(x_i) = d(x_0) + i$.
[On prend $y = F$, $x = \emptyset$ …]

\emph{Cor. 7} Soit $\Phi$ un polyèdre combinatoire strict, $d$ une fonction
dimension sur $\Phi$ (cf cor. 5) Soient $x, y \in \Phi$, $x < y$,
$d(y) = d(x) + 2$. Alors l'ens. des $z \in \Phi$ tels que $x < z < y$ est de
cardinal 2.

$\Phi_{xy}$ est un polyèdre combinatoire strict de dim $(-1) + 2 = 1$, donc
provient d'un polyèdre convexe de dim 1, i.e. d'un segment. Celui-ci a
exactement 2 extrémités !

\section{Cubes}

\note{« Cubes », à l'encre, est le seul mot de sa main sur la couverture de
papier brun qu'est la page 15 ; elle porte aussi, au crayon et entourée, la
mention « (13 p.) »}

\page{17}

Caractérisation intrinsèque d'un \uncertain{$n$-cube}, déterminée par un ens.
\[
  \Phi \subset \mathrm{Espaff}(E)
\]
\note{la notation est lue « $\mathrm{Espaff}$ » avec doute}
de sous-espaces affines d'un espace affine sur $k$ :

\begin{itemize}
  \item[a)] $\Phi$ \struck{a un plus petit et un plus grand élément, et}
    contient $\emptyset$ et $E$ (qui sont donc un plus petit et un plus grand
    élément), et dans $\Phi^{**} = \Phi - \{ 0_\Phi, 1_\Phi \}$ ($\Phi$ privé
    de son plus petit et plus grand élément) tout élément est intersection
    d'éléments maximaux (pour la rel. d'ordre).\note{l'exposant de
    $\Phi^{**}$ est fait de deux petits signes, lus $**$}
  \item[b)] Pour tout élément maximal $F$ de $\Phi^{**}$, $\exists !$ él.
    maximal $F'$ tel que $\operatorname{Inf}_{\Phi^{**}}(F, F') = \emptyset$
    [i.e. $(F, F')$ non minoré dans $\Phi^{**}$ i.e. $F \cap F'$ ne contient
    aucun él. de $\Phi$ différent de $\emptyset$] [Ainsi l'ens. $B$ des él.
    maximaux de $\Phi^{**}$ est muni d'une involution : $F \mapsto F'$ ss pt
    fixe …]
  \item[c)] Si \struck{:} à \uncertain{tout} $V \in \Phi^{**}$
    \struck{$V \neq \emptyset$} on associe l'ens. \struck{\ill{}} $B_V$ des
    $b \in B$ qui le majorent, on trouve une partie $A$ de $B$ telle que
    $b \in A \Rightarrow \sigma b \notin A$ (i.e. une section partielle
    \uncertain{définie} de $B$ sur $I = B/\sigma$). Ceci posé, l'application
    \[
      \Phi \setminus \{\emptyset\} \longrightarrow \Gamma_{\mathrm{part}}(B/I)
    \]
    ainsi définie, qui est injective par b), est bijection.
    \marginal{NB a), b), c) signifient que l'ens. ordonné $\Phi$ est un cube
    combinatoire}
  \item[d)] Si $n = \dim \mathrm{comb}\, \Phi$, on a \struck{$E$}
    $n =$ \struck{\ill{}} $\dim E = \frac{1}{2} \operatorname{card} B$ et
    $\forall V \in \Phi$, $\dim_\Phi V = \dim_E V$, \ill{}
    [a) : dit implique que \uncertain{la} $\operatorname{Inf}$
    \struck{\uncertain{dans} $\Phi$} \uncertain{d'une} partie minorée de
    $\Phi \setminus \{0\}$ est aussi son intersection (en tant
    \uncertain{qu'ens.} de points de $E$) et que les hyperplans $\in \Phi$
    \ill{}. Mais cela n'implique pas que le
\end{itemize}
\note{la fin de la page, écrite serrée et en partie en surcharge, n'est lue
que par fragments}

\page{18}

$\operatorname{Inf}$ d'une partie de $\Phi$ (dans $\Phi$) est $\emptyset$,
i.e. \uncertain{si} son intersection ne contient pas d'él. de
$\Phi \setminus \{\emptyset\}$, que cette intersection soit vide. Ex : un
\uncertain{quadrilatère} qui n'est pas un parallélogramme\note{suit un
petit croquis d'un quadrilatère dont deux côtés se coupent hors de lui}.
L'axiome qui suit assure qu'il en est pourtant ainsi …]

\begin{itemize}
  \item[e)] Si $b$, $b'$ sont des él. maximaux de \uncertain{$\Phi^{**}$}
    (i.e. $b, b' \in B$) dont l'inf est $\emptyset$ (i.e. $b' = \sigma b$),
    alors \struck{\ill{}} $b \cap b' = \emptyset$ (i.e., \struck{\ill{}} par
    d), \struck{\ill{}} $b$, $b'$ sont des \uncertain{parallèles} i.e.
    hyperplans affines parallèles …).
\end{itemize}

Dans l'espace des translations $T$ de $V$, les \struck{\ill{}} espaces
directeurs $\mathbb{H}_i$ ($i \in I$) des \struck{\ill{}} 2
\struck{sous-espaces \ill{}} hyperplans \ill{} $b \in B$ sont lin. indép.,
et correspondent donc à une décomposition
\[
  T \simeq \prod T_i \qquad (T_i = T/\mathbb{H}_i)
\]
de $T$ en produit de droites. La donnée du $T$-torseur $E$ revient à la
donnée d'une famille $(E_i)_{i \in I}$ de $T_i$-torseurs,
$E_i \simeq E/\mathbb{H}_i$. De plus, dans chaque $E_i$ il y a
\uncertain{une partie de} deux pts distincts, \struck{\ill{}} $B_i$,
(isom. à la fibre de $B/I$ en $i \in I$). Or on vérifie que la donnée
\struck{d'une} (pour $i$\note{la lettre rendue $\mathbb{H}$ est un $H$
barré de traits verticaux}

\page{19}

fixé) d'un $T_i$, $E_i$, $B_i \subset E_i$ équivaut en fait à la donnée de
$B_i$ (comme torseur sous $\pm 1$ !) en prenant
\[
  T_i = \operatorname{Ker}(k^{B_i} \xrightarrow{\varepsilon} k), \quad
  E_i = \varepsilon^{-1}(1)
\]
Donc la donnée \struck{des} d'une famille $(T_i, E_i, B_i)_{i \in I}$
($\operatorname{card} I = n$) équivaut à la donnée d'une \ill{} de
$n$-cube. Montrons que cette donnée permet de reconstruire le $n$-cube,
comme le produit de la famille (indexée par $I$) de 1-cubes …

On a aussi
\[
  T = \operatorname{Ker}(k^B \xrightarrow{\mathrm{tr}} k^I), \qquad
  E = \mathrm{tr}^{-1}(1)
\]
\note{« trace » est écrit au-dessus de la flèche, relié à « tr »}

Si $J \subset I$, on considère

\end{document}
